\documentclass[11pt]{article}
\usepackage[a4paper,margin=23mm]{geometry}
\usepackage{amsmath,amssymb,xcolor,hyperref,fancyhdr}
\usepackage{fontspec}
\setmainfont{Times New Roman}
\setmonofont{Times New Roman}
\usepackage[italic]{mathastext}
\setlength{\emergencystretch}{3em}
\setlength{\parskip}{7pt}\setlength{\parindent}{0pt}
\pagestyle{fancy}\fancyhf{}\setlength{\headheight}{14pt}
\lhead{FlowMLLab}\rhead{Reading and experiment guide}\cfoot{\thepage}
\begin{document}
\noindent\textbf{Week 4.2 | Reading the evidence}\par\medskip\hrule\medskip
\section*{Separate surrogate error from changing the CFD labels}
The low-fidelity Stokes field omits convection; the learned correction estimates
what must be added to approximate Navier\textendash{}Stokes on the specified grid. The
surrogate-to-label error and the coarse-to-fine CFD difference answer different
questions. A 0.1\% surrogate error against a coarse field can coexist with a 20\%
change in that field under refinement.

Pressure is reconstructed from momentum gradients with an aligned gauge. The
reported gradient residual measures compatibility of that reconstruction with
the supplied gradients; it is not a direct pressure-error norm and need not rank
models in the same order. Read velocity, pressure and compatibility together.
Optional retraining writes a new student directory. The retained figures remain
historical evidence; they do not become figures of your new fit.
\paragraph{Check your understanding.}
Identify the input information, the reference, the metric denominator and the
evaluation split in one worked notebook comparison. Predict which conclusion
would change if any one of those four changed, then check your prediction.
\end{document}
